% ==================== 4.1 引言 ====================
\section{引言}
\label{sec:ch4_intro}

巡视器动力学数字孪生建模是实现月面巡视器自主行走的关键技术之一。准确的动力学模型能够预测巡视器在月面环境下的运动行为，为路径规划、避障决策和运动控制提供支撑。本章针对月面巡视器的动力学特性，系统研究轮-壤交互力学理论和多体动力学建模方法，构建巡视器动力学数字孪生模型。

月面巡视器动力学建模面临以下挑战：

\textbf{（1）轮-壤相互作用复杂}

月面巡视器通常采用轮式移动方式，车轮与月壤的相互作用涉及复杂的接触力学、散体力学问题。月壤的低内聚力、高内摩擦角特性导致轮-壤相互作用表现出明显的滑移和沉陷行为，传统的刚性轮-刚性地面假设不再适用。

\textbf{（2）低重力环境影响}

月球重力仅为地球的1/6（约1.62 m/s²），低重力环境下巡视器的运动特性与地球环境显著不同。车轮接地压力减小，牵引能力下降，姿态稳定性变差。

\textbf{（3）多体动力学耦合}

巡视器由车身、悬架、车轮等多个刚体组成，各部件之间存在运动耦合。悬架的弹性变形、车轮的转向运动等因素对整体动力学行为有重要影响。

\textbf{（4）实时性要求高}

数字孪生系统要求动力学模型能够实时响应，这对模型复杂度和计算效率提出了约束。需要在模型精度和计算效率之间寻求平衡。

针对上述挑战，本章首先研究轮-壤交互力学理论，建立车轮与月壤相互作用的力学模型；然后基于多体动力学理论，构建巡视器整体动力学模型；最后通过仿真实验验证模型的有效性。
